\documentclass[12pt,a4paper]{article}

% ---------- Packages ----------
\usepackage[utf8]{inputenc}
\usepackage[T1]{fontenc}
\usepackage[french]{babel}
\usepackage{lmodern}
\usepackage[a4paper,margin=2.5cm]{geometry}
\usepackage{microtype}
\usepackage{setspace}
\usepackage{graphicx}
\usepackage{booktabs}
\usepackage{array}
\usepackage{xcolor}
\usepackage{listings}
\usepackage{pgfplots}
\usepackage{fancyhdr}
\usepackage{amsmath}
\usepackage[hidelinks]{hyperref}
\pgfplotsset{compat=1.16}

\setstretch{1.15}

% ---------- Style du code ----------
\definecolor{cmt}{RGB}{0,128,0}
\definecolor{kw}{RGB}{0,0,190}
\definecolor{bg}{RGB}{248,248,248}
\lstset{
  language=C++, basicstyle=\ttfamily\small, keywordstyle=\color{kw}\bfseries,
  commentstyle=\color{cmt}\itshape, stringstyle=\color{red!60!black},
  backgroundcolor=\color{bg}, frame=single, rulecolor=\color{black!25},
  numbers=left, numberstyle=\tiny\color{gray}, breaklines=true,
  showstringspaces=false, tabsize=4, captionpos=b, xleftmargin=1.5em,
  literate={é}{{\'e}}1 {è}{{\`e}}1 {à}{{\`a}}1
}

% ---------- En-tête / pied de page ----------
\pagestyle{fancy}
\fancyhf{}
\lhead{TP3 -- Recherche d'un mot dans un fichier}
\rhead{ILISI 2 -- 2026}
\cfoot{\thepage}
\renewcommand{\headrulewidth}{0.4pt}

\begin{document}

% ====================== PAGE DE GARDE ======================
\begin{titlepage}
  \centering
  \vspace*{1cm}
  {\large \textsc{Cycle d'ingénieur -- ILISI 2}\par}
  \vspace{0.3cm}
  {\large Année universitaire 2025--2026\par}
  \vspace{2cm}
  \rule{\linewidth}{1.2pt}\\[0.5cm]
  {\Huge \bfseries Travaux Pratiques n°3\par}
  \vspace{0.4cm}
  {\LARGE Recherche d'un mot dans un fichier texte\par}
  \vspace{0.3cm}
  {\Large Implémentation en C++ et \emph{benchmarking}\\ avec \texttt{<chrono>} et la STL\par}
  \rule{\linewidth}{1.2pt}\\[2cm]
  {\large Module : \textbf{C++ Avancé et Virtualisation}\par}
  \vspace{2cm}
  {\large Réalisé par :\par}
  {\Large \textbf{YOUSSEF ERRAMI}\par}
  \vfill
  {\large \today\par}
\end{titlepage}

\tableofcontents
\newpage

% ====================== 1. INTRODUCTION ======================
\section{Introduction}

La recherche d'une sous-chaîne dans un texte est une opération fondamentale
en informatique : éditeurs, moteurs de recherche, outils comme \texttt{grep}.
Ce TP consiste à écrire un programme C++ qui lit un fichier \texttt{.txt}, demande
un mot à l'utilisateur et indique s'il existe dans le fichier, sans tenir compte de la casse.

Deux implémentations de même logique sont comparées :
\begin{enumerate}
  \item une version \textbf{STL} reposant sur \texttt{std::string} et \texttt{std::string::find} ;
  \item une version \textbf{\emph{from scratch}}, où l'algorithme de recherche naïf est écrit à la main.
\end{enumerate}
Les temps d'exécution sont mesurés avec la bibliothèque \texttt{<chrono>} et exploités
statistiquement à l'aide de \texttt{std::vector}, \texttt{std::sort} et \texttt{std::accumulate}.

\paragraph{Objectifs.}
(i) Implémenter les deux versions ; (ii) mesurer leurs temps de lecture et de recherche ;
(iii) comparer leur comportement selon la taille du fichier et la position du mot ;
(iv) interpréter les écarts observés.

% ====================== 2. ENVIRONNEMENT ======================
\section{Environnement et organisation du projet}

\begin{table}[h]
\centering
\begin{tabular}{ll}
\toprule
\textbf{Élément} & \textbf{Valeur} \\
\midrule
Système & Windows 11 Pro \\
Compilateur & g++ 16.2.0 (MinGW-w64, UCRT) \\
Standard & C++17 \\
Options & \texttt{-O2 -Wall -Wextra -static} \\
Horloge & \texttt{std::chrono::steady\_clock} \\
\bottomrule
\end{tabular}
\caption{Configuration de test.}
\end{table}

\begin{table}[h]
\centering
\begin{tabular}{p{5.6cm}p{8.4cm}}
\toprule
\textbf{Fichier} & \textbf{Rôle} \\
\midrule
\texttt{recherche\_fichier\_string.cpp} & Version STL (\texttt{std::string::find}) \\
\texttt{recherche\_fichier\_scratch.cpp} & Version \emph{from scratch} (algorithme naïf) \\
\texttt{benchmark.cpp} & Benchmark comparatif des deux versions (sortie CSV) \\
\texttt{texte.txt} & Fichier de test (130 lignes, 6\,711 octets) \\
\texttt{resultats\_benchmark.csv} & Résultats bruts des mesures \\
\bottomrule
\end{tabular}
\caption{Organisation du projet.}
\end{table}

Chaque programme est structuré en sections : utilitaires, lecture, recherche,
affichage, benchmark, puis \texttt{main}. La logique d'origine n'a pas été modifiée.

% ====================== 3. CONCEPTION ======================
\section{Conception}

\subsection{Version STL}
Le fichier est lu en une seule opération par \texttt{rdbuf()} dans un
\texttt{std::stringstream}. Le texte et le mot sont ensuite convertis en minuscules
avec \texttt{std::transform}, puis \texttt{std::string::find} renvoie la position de la
première occurrence ou \texttt{std::string::npos}.

\begin{lstlisting}[caption={Recherche de la version STL.}]
size_t rechercher(const std::string& contenu, const std::string& mot) {
    std::string contenuLower = toLower(contenu);
    std::string motLower     = toLower(mot);
    return contenuLower.find(motLower);
}
\end{lstlisting}

\subsection{Version \emph{from scratch}}
Le fichier est lu caractère par caractère avec \texttt{get()}. La recherche teste
chaque position $i$ du texte et compare le motif caractère par caractère, en comparant
\texttt{tolower} des deux côtés : aucune copie du texte n'est nécessaire.

\begin{lstlisting}[caption={Recherche de la version \emph{from scratch}.}]
int rechercher(const std::string& texte, const std::string& motif) {
    if (motif.empty() || motif.size() > texte.size()) return -1;
    int n = (int)texte.size(), m = (int)motif.size();
    for (int i = 0; i <= n - m; ++i) {
        int j = 0;
        while (j < m && std::tolower((unsigned char)texte[i + j]) ==
                        std::tolower((unsigned char)motif[j]))
            ++j;
        if (j == m) return i;   // motif trouvé en i
    }
    return -1;
}
\end{lstlisting}

\subsection{Complexité théorique}
Soit $n$ la taille du texte et $m$ celle du motif.

\begin{table}[h]
\centering
\begin{tabular}{lcc}
\toprule
 & \textbf{Pire cas} & \textbf{Cas favorable} \\
\midrule
Version STL (copie + \texttt{find}) & $\Theta(n)$ systématique (copie) & $\Theta(n)$ \\
Version \emph{from scratch} & $O(n\cdot m)$ & $O(m)$ (mot au début) \\
\bottomrule
\end{tabular}
\caption{Complexités des deux approches.}
\end{table}

La version STL paie toujours le coût de la mise en minuscules de tout le texte, même
si le mot se trouve au début. La version naïve peut s'arrêter dès la première occurrence.

% ====================== 4. MÉTHODOLOGIE ======================
\section{Méthodologie de mesure}

\subsection{Chronométrage avec \texttt{<chrono>}}
Un modèle de fonction mesure la durée d'un appelable en microsecondes avec une horloge
monotone (\texttt{steady\_clock}), non affectée par les changements d'heure système.

\begin{lstlisting}[caption={Fonction de chronométrage.}]
template <typename Fonction>
double mesurerUs(Fonction f) {
    auto debut = Horloge::now();
    f();
    auto fin = Horloge::now();
    return std::chrono::duration<double, std::micro>(fin - debut).count();
}
\end{lstlisting}

\subsection{Statistiques avec la STL}
Chaque mesure est répétée $N$ fois et les durées sont rangées dans un
\texttt{std::vector<double>}. Après \texttt{std::sort}, on obtient le minimum, le
maximum et la médiane ; \texttt{std::accumulate} donne la moyenne et l'écart-type.
La variable \texttt{volatile} évite que l'optimiseur supprime l'appel mesuré.

\begin{lstlisting}[caption={Collecte et analyse des temps.}]
std::vector<double> temps;
for (int i = 0; i < repetitions; ++i)
    temps.push_back(mesurerUs([&]{ sink = rechercher(contenu, mot); }));
std::sort(temps.begin(), temps.end());
double moyenne = std::accumulate(temps.begin(), temps.end(), 0.0) / temps.size();
double mediane = temps[temps.size() / 2];
\end{lstlisting}

\subsection{Protocole expérimental}
\begin{itemize}
  \item \textbf{Tailles :} le fichier de base (6\,711 octets) est répliqué $\times 1$, $\times 10$,
        $\times 100$ et $\times 1000$ (jusqu'à environ 6,7 Mo).
  \item \textbf{Mots :} \og Bonjour \fg{} (début du texte), \og Fin du fichier de test \fg{}
        (fin de la première copie) et \og zzzxyz \fg{} (absent : parcours complet).
  \item \textbf{Répétitions :} 200 par couple (version, mot, taille).
\end{itemize}

% ====================== 5. RÉSULTATS ======================
\section{Résultats}

\subsection{Programmes individuels}
Sur le fichier de base avec le mot \og compteur \fg, les deux programmes trouvent
l'occurrence à la position 118. Pour la version \emph{from scratch}, la recherche est
mesurée à environ 0,8~$\mu$s (médiane sur 1000 répétitions), la lecture complète
(caractère par caractère) à environ 260~$\mu$s.

\subsection{Benchmark comparatif}
Les temps médians (en $\mu$s) sont donnés par la table~\ref{tab:resultats}.

\begin{table}[h]
\centering
\small
\begin{tabular}{rl rr}
\toprule
\textbf{Taille (octets)} & \textbf{Mot} & \textbf{STL} & \textbf{Scratch} \\
\midrule
6\,711     & début  & 17,10     & 0,10 \\
           & fin    & 17,80     & 33,10 \\
           & absent & 17,50     & 32,50 \\
\midrule
67\,110    & début  & 200,10    & 0,10 \\
           & fin    & 198,30    & 32,20 \\
           & absent & 194,00    & 340,60 \\
\midrule
671\,100   & début  & 1\,675,00 & 0,10 \\
           & fin    & 1\,682,10 & 31,40 \\
           & absent & 1\,714,80 & 3\,234,60 \\
\midrule
6\,711\,000 & début  & 17\,775,10 & 0,10 \\
            & fin    & 18\,448,90 & 32,30 \\
            & absent & 23\,804,70 & 49\,497,40 \\
\bottomrule
\end{tabular}
\caption{Temps médians de recherche ($\mu$s), 200 répétitions.}
\label{tab:resultats}
\end{table}

\begin{figure}[h]
\centering
\begin{tikzpicture}
\begin{loglogaxis}[
  width=0.88\linewidth, height=7.5cm,
  xlabel={Taille du texte (octets)}, ylabel={Temps médian ($\mu$s)},
  grid=major, legend pos=north west, legend cell align=left,
  legend style={font=\small}
]
\addplot[blue, mark=*, thick] coordinates {(6711,17.5)(67110,194)(671100,1714.8)(6711000,23804.7)};
\addlegendentry{STL -- mot absent}
\addplot[red, mark=square*, thick] coordinates {(6711,32.5)(67110,340.6)(671100,3234.6)(6711000,49497.4)};
\addlegendentry{Scratch -- mot absent}
\addplot[blue, dashed, mark=o, thick] coordinates {(6711,17.1)(67110,200.1)(671100,1675)(6711000,17775.1)};
\addlegendentry{STL -- mot au début}
\addplot[red, dashed, mark=square, thick] coordinates {(6711,0.1)(67110,0.1)(671100,0.1)(6711000,0.1)};
\addlegendentry{Scratch -- mot au début}
\end{loglogaxis}
\end{tikzpicture}
\caption{Temps de recherche en fonction de la taille du texte (échelle log-log).}
\label{fig:perf}
\end{figure}

% ====================== 6. ANALYSE ======================
\section{Analyse et discussion}

\paragraph{Croissance linéaire.}
Dans le pire cas (mot absent), les deux versions suivent une droite de pente 1 en
échelle log-log : le temps est multiplié par environ 10 quand la taille l'est aussi,
ce qui confirme une complexité $O(n)$ pour ce type de texte.

\paragraph{Mot en début de texte.}
La version \emph{from scratch} s'arrête dès la première occurrence : $\approx 0{,}1~\mu$s
quelle que soit la taille. La version STL reste proportionnelle à $n$, car elle recopie
et convertit \emph{tout} le texte en minuscules avant de chercher. Dans ce cas, l'écart
atteint plus de cinq ordres de grandeur sur 6,7~Mo.

\paragraph{Mot absent.}
La version STL est environ 1,8 à 2,1 fois plus rapide : \texttt{find} s'appuie sur des
primitives optimisées, tandis que la version naïve appelle \texttt{tolower} deux fois à chaque
comparaison. Le coût dominant de la version STL est la copie avec conversion, pas la recherche.

\paragraph{Variabilité.}
L'écart-type devient notable sur les gros fichiers (par ex. 6\,706~$\mu$s pour la version
STL sur 6,7~Mo, mot absent), du fait des effets de cache, des allocations mémoire et de
l'ordonnancement du système. C'est pourquoi la médiane est retenue plutôt que la moyenne.

% ====================== 7. LIMITES ======================
\section{Limites et améliorations possibles}

\begin{itemize}
  \item La recherche insensible à la casse ne traite que l'ASCII ; les caractères accentués
        UTF-8 nécessitent une bibliothèque Unicode.
  \item Le mot \og fin \fg{} est trouvé dans la première copie du texte répliqué, d'où des
        temps constants pour la version \emph{from scratch}.
  \item Les mesures dépendent de la machine ; elles servent à comparer les versions
        entre elles, pas à fournir des valeurs absolues.
  \item \textbf{Améliorations :} éviter la copie dans la version STL en comparant avec
        \texttt{std::search} et un prédicat insensible à la casse ; utiliser un algorithme
        sous-linéaire (Boyer--Moore--Horspool, \texttt{std::boyer\_moore\_searcher}) ;
        lire le fichier par blocs.
\end{itemize}

% ====================== 8. CONCLUSION ======================
\section{Conclusion}

Ce TP a permis d'implémenter la recherche d'un mot dans un fichier de deux manières
et de les comparer quantitativement avec \texttt{<chrono>} et les conteneurs de la STL.
Les résultats montrent qu'aucune version n'est meilleure dans tous les cas : la version
\emph{from scratch} est très rapide quand le mot est proche du début, mais elle est
environ deux fois plus lente que la STL dans le pire cas ; la version STL a un coût
constant en $O(n)$ lié à la copie du texte. Le benchmarking est donc indispensable pour
valider un choix d'implémentation, à condition de répéter les mesures et de s'appuyer
sur des statistiques robustes comme la médiane.

% ====================== ANNEXES ======================
\appendix
\section{Compilation et exécution}

\begin{lstlisting}[language=bash,numbers=none,caption={Commandes utilisées.}]
g++ -std=c++17 -O2 -Wall -Wextra -static recherche_fichier_string.cpp  -o recherche_fichier_string.exe
g++ -std=c++17 -O2 -Wall -Wextra -static recherche_fichier_scratch.cpp -o recherche_fichier_scratch.exe
g++ -std=c++17 -O2 -Wall -Wextra -static benchmark.cpp                 -o benchmark.exe

recherche_fichier_string.exe  texte.txt
recherche_fichier_scratch.exe texte.txt
benchmark.exe texte.txt > resultats_benchmark.csv
\end{lstlisting}

L'option \texttt{-static} évite le chargement d'une \texttt{libstdc++-6.dll} incompatible
présente dans le \texttt{PATH} (celle de Git for Windows), qui provoquait un plantage au démarrage.

\end{document}
